% Adapted from the author's linear-algebra book as a full-width comparison.
\begin{tikzpicture}[
  x=1mm,y=1mm,font=\footnotesize,line width=0.45pt,
  every node/.style={inner sep=2pt},
  map arrow/.style={line width=0.35pt,-{Latex[length=1.6mm,width=1.1mm]},
    shorten <=2.5pt,shorten >=2.5pt}
]
  \foreach \origin/\paneltitle in {
    0/{Injective, not surjective},
    52/{Surjective, not injective},
    104/{Bijective}
  } {
    \begin{scope}[xshift=\origin mm]
      \node at (21.5,28) {\paneltitle};
      \draw[fill=black!5] (9,0) ellipse[x radius=7mm,y radius=17mm];
      \draw[fill=black!5] (34,0) ellipse[x radius=7mm,y radius=17mm];
      \node at (9,22) {\(A\)};
      \node at (34,22) {\(B\)};
    \end{scope}
  }
  \begin{scope}
    \foreach \y in {9,0,-9} {
      \fill (9,\y) circle[radius=1.6pt];
    }
    \foreach \y in {9,3,-3,-9} {
      \fill (34,\y) circle[radius=1.6pt];
    }
    \draw[map arrow] (9,9) -- (34,9);
    \draw[map arrow] (9,0) -- (34,3);
    \draw[map arrow] (9,-9) -- (34,-3);
  \end{scope}
  \begin{scope}[xshift=52mm]
    \foreach \y in {9,3,-3,-9} {
      \fill (9,\y) circle[radius=1.6pt];
    }
    \foreach \y in {9,0,-9} {
      \fill (34,\y) circle[radius=1.6pt];
    }
    \draw[map arrow] (9,9) -- (34,9);
    \draw[map arrow] (9,3) -- (34,0);
    \draw[map arrow] (9,-3) -- (34,0);
    \draw[map arrow] (9,-9) -- (34,-9);
  \end{scope}
  \begin{scope}[xshift=104mm]
    \foreach \y in {9,0,-9} {
      \fill (9,\y) circle[radius=1.6pt];
      \fill (34,\y) circle[radius=1.6pt];
      \draw[map arrow] (9,\y) -- (34,\y);
    }
  \end{scope}
\end{tikzpicture}
